\begin{ex}
Đơn vị đo diện tích xăng-ti-mét vuông viết tắt là gì? (Viết số: $\text{cm}^2$ có số mũ là bao nhiêu?).
\shortans{2}
\loigiai{
Số mũ là 2.
}
\end{ex}

\begin{ex}
Một hình vuông có cạnh dài 1 cm. Diện tích của hình vuông đó là bao nhiêu xăng-ti-mét vuông?
\shortans{1}
\loigiai{
$1 \times 1 = 1\text{ cm}^2$.
}
\end{ex}

\begin{ex}
Một hình chữ nhật được ghép bởi 6 ô vuông nhỏ, mỗi ô vuông có diện tích $1\text{ cm}^2$. Diện tích hình chữ nhật đó là bao nhiêu xăng-ti-mét vuông?
\shortans{6}
\loigiai{
$6 \times 1 = 6\text{ cm}^2$.
}
\end{ex}

\begin{ex}
Một tấm bìa được tạo bởi 15 ô vuông nhỏ diện tích $1\text{ cm}^2$. Diện tích tấm bìa là bao nhiêu xăng-ti-mét vuông?
\shortans{15}
\loigiai{
$15\text{ cm}^2$.
}
\end{ex}

\begin{ex}
Tính: $25\text{ cm}^2 + 35\text{ cm}^2 = \dots\text{ cm}^2$.
\shortans{60}
\loigiai{
$25 + 35 = 60$.
}
\end{ex}

\begin{ex}
Tính: $90\text{ cm}^2 - 42\text{ cm}^2 = \dots\text{ cm}^2$.
\shortans{48}
\loigiai{
$90 - 42 = 48$.
}
\end{ex}

\begin{ex}
Tính: $15\text{ cm}^2 \times 4 = \dots\text{ cm}^2$.
\shortans{60}
\loigiai{
$15 \times 4 = 60$.
}
\end{ex}

\begin{ex}
Tính: $84\text{ cm}^2 : 4 = \dots\text{ cm}^2$.
\shortans{21}
\loigiai{
$84 : 4 = 21$.
}
\end{ex}

\begin{ex}
Một miếng bìa diện tích $45\text{ cm}^2$, người ta cắt đi $18\text{ cm}^2$. Miếng bìa còn lại bao nhiêu xăng-ti-mét vuông?
\shortans{27}
\loigiai{
$45 - 18 = 27\text{ cm}^2$.
}
\end{ex}

\begin{ex}
Một hình gồm 8 ô vuông, mỗi ô vuông có diện tích $2\text{ cm}^2$. Diện tích cả hình là bao nhiêu xăng-ti-mét vuông?
\shortans{16}
\loigiai{
$8 \times 2 = 16\text{ cm}^2$.
}
\end{ex}

\begin{ex}
Tính: $30\text{ cm}^2 + 20\text{ cm}^2 \times 2$.
\shortans{70}
\loigiai{
$30 + 40 = 70$.
}
\end{ex}

\begin{ex}
Tính: $(50\text{ cm}^2 - 20\text{ cm}^2) : 3$.
\shortans{10}
\loigiai{
$30 : 3 = 10$.
}
\end{ex}

\begin{ex}
Một hình chữ nhật có diện tích $24\text{ cm}^2$ được chia thành 4 phần bằng nhau. Mỗi phần có diện tích bao nhiêu xăng-ti-mét vuông?
\shortans{6}
\loigiai{
$24 : 4 = 6\text{ cm}^2$.
}
\end{ex}

\begin{ex}
Một hình tam giác có diện tích $18\text{ cm}^2$, một hình tứ giác có diện tích $25\text{ cm}^2$. Tổng diện tích hai hình là bao nhiêu xăng-ti-mét vuông?
\shortans{43}
\loigiai{
$18 + 25 = 43\text{ cm}^2$.
}
\end{ex}

\begin{ex}
Hình A có diện tích $36\text{ cm}^2$, hình B có diện tích $12\text{ cm}^2$. Diện tích hình A gấp mấy lần diện tích hình B?
\shortans{3}
\loigiai{
$36 : 12 = 3$ lần.
}
\end{ex}

\begin{ex}
Hình M có diện tích $40\text{ cm}^2$, hình N có diện tích kém hình M là $15\text{ cm}^2$. Diện tích hình N là bao nhiêu xăng-ti-mét vuông?
\shortans{25}
\loigiai{
$40 - 15 = 25\text{ cm}^2$.
}
\end{ex}

\begin{ex}
Tính: $12\text{ cm}^2 \times 5 = \dots\text{ cm}^2$.
\shortans{60}
\loigiai{
$12 \times 5 = 60$.
}
\end{ex}

\begin{ex}
Tính: $100\text{ cm}^2 : 5 = \dots\text{ cm}^2$.
\shortans{20}
\loigiai{
$100 : 5 = 20$.
}
\end{ex}

\begin{ex}
Một tờ giấy màu đỏ diện tích $50\text{ cm}^2$, tờ giấy màu xanh diện tích $70\text{ cm}^2$. Cả hai tờ giấy có diện tích bao nhiêu xăng-ti-mét vuông?
\shortans{120}
\loigiai{
$50 + 70 = 120\text{ cm}^2$.
}
\end{ex}

\begin{ex}
Một con tem hình chữ nhật có diện tích $6\text{ cm}^2$. Hỏi 5 con tem như thế có tổng diện tích bao nhiêu xăng-ti-mét vuông?
\shortans{30}
\loigiai{
$6 \times 5 = 30\text{ cm}^2$.
}
\end{ex}

\begin{ex}
Tấm bìa hình vuông gồm 16 ô vuông nhỏ mỗi ô diện tích $1\text{ cm}^2$. Hỏi diện tích tấm bìa là bao nhiêu xăng-ti-mét vuông?
\shortans{16}
\loigiai{
$16 \times 1 = 16\text{ cm}^2$.
}
\end{ex}

\begin{ex}
Tính: $14\text{ cm}^2 + 36\text{ cm}^2 : 4$.
\shortans{23}
\loigiai{
$14 + 9 = 23$.
}
\end{ex}

\begin{ex}
Tính: $60\text{ cm}^2 - 15\text{ cm}^2 \times 2$.
\shortans{30}
\loigiai{
$60 - 30 = 30$.
}
\end{ex}

\begin{ex}
Một miếng dán sticker có diện tích $8\text{ cm}^2$. Một trang vở dán được 9 miếng như vậy. Tổng diện tích các miếng sticker là bao nhiêu xăng-ti-mét vuông?
\shortans{72}
\loigiai{
$8 \times 9 = 72\text{ cm}^2$.
}
\end{ex}

\begin{ex}
Hình vuông cạnh 1 cm thì chu vi là 4 cm, còn diện tích là bao nhiêu xăng-ti-mét vuông?
\shortans{1}
\loigiai{
$1 \times 1 = 1\text{ cm}^2$.
}
\end{ex}
